\section{Learning how far to edit}
The preceding bound motivates removing the inherited displacement cap.
We train two response models on five fractions of a causal candidate path
from the CV reference, using the same training draws and
\AmplitudeStudyUpdates{} total updates. For a fixed path direction and
label mask, the Euclidean error is convex in the edit fraction.
The displacement debit on unknown future steps is linear for nonnegative
fractions. These properties motivate monotone slopes for concave lower
gain and convex harm and signed change. The second model predicts the
same quantities directly at the five fractions. Both predict the
reference masses separately. Neither sees future targets or masks at
inference. Training and calibration use analytic candidate directions;
the held development fold adds source-excluded neural directions.

Across the nine held development sources, the shape constraint reduces
normalized dual-gain MAE by \AmplitudeStudyGainError{} and signed easy
MAE by \AmplitudeStudySignedError{} (3,000 paired source resamples).
This better fit does not translate to better decisions. All effects below
are paired completion-lower differences in percentage points with 95\%
source-bootstrap intervals:
\begin{itemize}
  \item Shape versus direct independent ranking at the same edit
  counts in each query: \AmplitudeStudyCount{}.
  \item Shape query portfolio versus independent shape selection
  at the same query counts: \AmplitudeStudyPortfolio{}.
  \item Shape portfolio versus uniform amplitudes on identical rows
  and rays, conserving the sum of peak displacement per query:
  \AmplitudeStudyMass{}.
\end{itemize}
The shape portfolio makes \AmplitudeStudyEdits{} edits.
It passes \AmplitudeStudyRisk{}/72 original risk checks, versus
\AmplitudeStudyKnown{}/72 if missing-step uncertainty is ignored.
Its worst known signed easy change is \AmplitudeStudyWorst{}\%;
doubling the frozen edits instead gives \AmplitudeStudyBoostWorst{}\%.
Its lower contrasts are \AmplitudeStudyCV{} against CV and
\AmplitudeStudyOLS{} against OLS4. No calibration cell qualifies
(\AmplitudeStudyQual{}/9); neither a larger amplitude nor a more
accurate response model establishes an eligible policy.
These are exposed development results, not external confirmation.
